\chapter{System Design}
\label{chap:design}

\section{Introduction}

This chapter presents the architecture of CANvisionNative: the overall layered architecture, the desktop client, the backend, the database, the \gls{rest} \gls{api}, the replay engine, the hardware layer, the detection pipeline, and the \gls{ai} explanation engine. It closes with the class and entity-relationship diagrams used to implement the design described here.

\section{Overall architecture}

CANvisionNative is built as two cooperating processes on the same machine: a Windows desktop client and a local Python backend, communicating over \gls{http} on \texttt{localhost}. Figure~\ref{fig:overall-arch} shows the layered view of the system.

\begin{figure}[H]
\centering
\begin{tikzpicture}[
    layer/.style={draw, rounded corners, minimum width=12cm, minimum height=1.3cm, align=center, font=\small},
    >=Stealth
]
\node[layer, fill=blue!8] (client) at (0,6) {\textbf{Desktop client} — WPF / MVVM (.NET Framework 4.8)\\ Views, ViewModels, VehicleDataService, PythonApiClient};
\node[layer, fill=orange!8] (api) at (0,4.2) {\textbf{REST API layer} — FastAPI (\texttt{backend/api/server.py})\\ replay, live, explain, offline-analyze, hardware, settings endpoints};
\node[layer, fill=green!8] (engine) at (0,2.4) {\textbf{Detection engine} — RealtimeEngine, ScoringEngine,\\ RuntimeModelStore (per-vehicle Isolation Forest + KMeans)};
\node[layer, fill=yellow!12] (ai) at (0,0.6) {\textbf{AI explanation engine} — rule-based reasoning,\\ context builder, LLM provider chain (Gemini $\to$ Ollama $\to$ template)};
\node[layer, fill=gray!10] (data) at (0,-1.2) {\textbf{Data layer} — DBC files, trained model store (.joblib),\\ SQLite database, JSON knowledge base, log files};

\draw[<->] (client) -- (api);
\draw[<->] (api) -- (engine);
\draw[<->] (engine) -- (ai);
\draw[<->] (api) -- (data);
\draw[<->] (engine) -- (data);
\draw[<->] (ai) -- (data);
\end{tikzpicture}
\caption{Overall layered architecture of CANvisionNative.}
\label{fig:overall-arch}
\end{figure}

\subsection{Frontend architecture}

The desktop client follows the \gls{mvvm} pattern. Each screen (Home, Dashboard, Telemetry, Diagnostics, Log Playback, Anomaly Intel, Settings) has a XAML view and a corresponding ViewModel in \texttt{ViewModels/SectionViewModels.cs}. Two services sit between the ViewModels and the backend, as shown in Figure~\ref{fig:frontend-arch}:

\begin{itemize}
    \item \textbf{PythonApiClient} — a thin \gls{http} client wrapping every backend endpoint (start/stop replay, get alerts, get explanation, get vehicle list, and so on).
    \item \textbf{VehicleDataService} — the shared, in-memory state of the client: current snapshot, alert history, replay runtime metrics. It owns a polling timer that periodically calls the backend (replay status, partial alerts, live signals) and raises events that the ViewModels subscribe to, so that every view reflects the same underlying state.
\end{itemize}

\begin{figure}[H]
\centering
\begin{tikzpicture}[
    box/.style={draw, rounded corners, minimum width=3.6cm, minimum height=0.9cm, align=center, font=\small},
    >=Stealth, node distance=1.1cm and 1.4cm
]
\node[box, fill=blue!8] (views) {XAML Views\\(Home, Dashboard, \ldots)};
\node[box, fill=orange!8, below=of views] (vms) {ViewModels\\(SectionViewModels.cs)};
\node[box, fill=green!8, below left=of vms] (vds) {VehicleDataService\\(shared state + poll timer)};
\node[box, fill=green!8, below right=of vms] (api) {PythonApiClient\\(HTTP wrapper)};
\node[box, fill=gray!10, below=1.6cm of vms] (backend) {Backend REST API};

\draw[<->] (views) -- (vms);
\draw[<->] (vms) -- (vds);
\draw[<->] (vms) -- (api);
\draw[<->] (vds) -- (backend);
\draw[<->] (api) -- (backend);
\end{tikzpicture}
\caption{Frontend architecture: MVVM views bound to ViewModels, backed by two shared services.}
\label{fig:frontend-arch}
\end{figure}

Navigation between sections is driven by a section catalog and a navigation service, so that adding a new screen does not require changes to existing screens.

\subsection{Backend architecture}

The backend is a single FastAPI application (\texttt{backend/api/server.py}) organized around five responsibilities:

\begin{enumerate}
    \item \textbf{Decoding} (\texttt{backend/decoding/}) — vehicle profile registry and \gls{dbc}-based signal decoding.
    \item \textbf{Data processing} (\texttt{backend/data\_processing/}) — log parsing and \gls{mf4}-to-\gls{csv} conversion.
    \item \textbf{Machine learning} (\texttt{backend/ml/}) — feature engineering, the runtime detection engine, training scripts, and the \gls{ai} explainer.
    \item \textbf{Hardware} (\texttt{backend/hardware/}) — the live \gls{can} adapter interface and session recorder.
    \item \textbf{Persistence} (\texttt{backend/db/}) — the SQLite session and alert store.
\end{enumerate}

Long-running work (a replay, an offline analysis) is executed in a background subprocess or thread so that the \gls{api} stays responsive; the client polls a status endpoint until the work completes.

\subsection{Database}

The backend keeps a small SQLite database (\texttt{backend/db/canvision.db}) for session bookkeeping and historical alerts, independent of the per-replay summary files used for the live client views. Figure~\ref{fig:er-diagram} shows its three tables.

\begin{figure}[H]
\centering
\begin{tikzpicture}[
    entity/.style={draw, rectangle, minimum width=4.6cm, align=left, font=\scriptsize},
    title/.style={font=\bfseries\small},
    >=Stealth
]
\node[entity] (sessions) at (0,0) {
    \underline{\textbf{sessions}}\\
    \underline{id}: TEXT\\
    started\_at: REAL\\
    ended\_at: REAL\\
    source: TEXT (live/replay/inject)\\
    frame\_count: INTEGER\\
    anomaly\_count: INTEGER
};
\node[entity] (frames) at (6.4,1.6) {
    \underline{\textbf{live\_frames}}\\
    \underline{id}: INTEGER (PK)\\
    session\_id: TEXT (FK)\\
    timestamp: REAL\\
    can\_id: INTEGER\\
    b0..b7: INTEGER\\
    anomaly: INTEGER\\
    user\_label: TEXT
};
\node[entity] (alerts) at (6.4,-1.6) {
    \underline{\textbf{alert\_history}}\\
    \underline{id}: INTEGER (PK)\\
    session\_id: TEXT (FK)\\
    timestamp: REAL\\
    can\_id: TEXT\\
    severity: TEXT\\
    score: REAL\\
    attack\_type, reason: TEXT\\
    temporal\_pattern: TEXT\\
    plausibility\_passed: INTEGER\\
    timing\_anomaly: INTEGER\\
    recommended\_response: TEXT
};
\draw[-{Latex[length=3mm]}] (sessions) -- (frames) node[midway,above,sloped,font=\tiny]{1 .. N};
\draw[-{Latex[length=3mm]}] (sessions) -- (alerts) node[midway,below,sloped,font=\tiny]{1 .. N};
\end{tikzpicture}
\caption{Entity-relationship diagram of the SQLite persistence layer.}
\label{fig:er-diagram}
\end{figure}

A \texttt{session} groups the frames and alerts produced by one live, replay, or injection run. \texttt{live\_frames} keeps every captured frame with its raw bytes and the engine's normal/anomalous classification, including an optional \texttt{user\_label} so a human reviewer can correct the label for future retraining. \texttt{alert\_history} keeps one row per raised alert, including the parsed detection-layer reason string and the recommended response, indexed on \texttt{session\_id} and \texttt{severity} for fast filtering.

\section{REST API}

The backend exposes more than 40 endpoints under a small number of prefixes: \texttt{/api/replay/*} (start, stop, status, alerts), \texttt{/api/live/*} (simulation and hardware capture), \texttt{/api/explain/\{index\}} (\gls{ai} explanation), \texttt{/api/offline/*} (batch analysis), \texttt{/api/vehicles} (profile registry), \texttt{/api/session/*} (recording), and \texttt{/api/settings}. The full endpoint list is given in Appendix~\ref{app:api}. All endpoints return \gls{json} and are consumed exclusively by the desktop client over \texttt{localhost}.

\section{Replay engine}
\label{sec:design-replay}

The replay engine (\texttt{backend/ml/runtime/replay\_runner.py}) is launched as a subprocess by the \texttt{/api/replay/start} endpoint. It reads the input file through a single, format-aware loader (Section~\ref{sec:impl-parser}) shared with the offline batch analyzer, so that every supported format enters the pipeline as the same canonical \texttt{(timestamp, can\_id, payload)} stream. Each frame is then passed, in order, to the same \texttt{RealtimeEngine} used by the live and simulated paths, guaranteeing that replay, live, and simulated detection behave identically for the same input. Section~\ref{sec:impl-replay} describes its implementation in detail, and Chapter~\ref{chap:validation} documents two defects that were found and fixed in this exact module.

\section{Detection pipeline}
\label{sec:design-pipeline}

Figure~\ref{fig:pipeline} shows the per-frame detection pipeline, common to replay, simulation, and live hardware capture.

\begin{figure}[H]
\centering
\begin{tikzpicture}[
    box/.style={draw, rounded corners, minimum width=3.1cm, minimum height=1.1cm, align=center, font=\small},
    >=Stealth, node distance=0.5cm
]
\node[box, fill=blue!8] (frame) {CAN frame\\(timestamp, id, bytes)};
\node[box, fill=orange!8, right=of frame] (feat) {Feature\\extraction};
\node[box, fill=green!8, above right=0.3cm and 1.2cm of feat] (timing) {Timing\\layer};
\node[box, fill=green!8, right=1.2cm of feat] (payload) {Payload\\layer};
\node[box, fill=green!8, below right=0.3cm and 1.2cm of feat] (ml) {ML layer\\(per-vehicle Isolation Forest)};
\node[box, fill=purple!10, right=2.6cm of payload] (fusion) {Fusion score\\+ persistence};
\node[box, fill=red!10, right=of fusion] (alert) {Alert +\\severity};

\draw[->] (frame) -- (feat);
\draw[->] (feat) -- (timing);
\draw[->] (feat) -- (payload);
\draw[->] (feat) -- (ml);
\draw[->] (timing) -- (fusion);
\draw[->] (payload) -- (fusion);
\draw[->] (ml) -- (fusion);
\draw[->] (fusion) -- (alert);
\end{tikzpicture}
\caption{Per-frame detection pipeline: feature extraction, four-layer scoring, fusion, and alerting.}
\label{fig:pipeline}
\end{figure}

Feature extraction computes, per \gls{can} identifier, a rolling set of statistics: inter-arrival time and its variance, message frequency, byte-level difference from the previous frame on the same identifier, and Shannon entropy of the identifier stream and of the payload. These features feed three of the four layers directly:

\begin{itemize}
    \item The \textbf{timing layer} flags both abnormally bursty timing and abnormally high message frequency (a pure timing-variance metric alone would miss a steady flood, since a perfectly periodic flood has almost no timing variance).
    \item The \textbf{payload layer} flags abrupt, out-of-pattern byte changes.
    \item The \textbf{ML layer} first classifies the vehicle's current operating state (parking, idle, driving) with a K-Means model, then scores the feature vector with the Isolation Forest trained for that state and that specific vehicle profile, falling back to a lightweight heuristic if no matching per-vehicle model is available.
    \item The \textbf{temporal-persistence layer} tracks, over a rolling window, how often the fused score has stayed above the high-severity threshold, so an isolated single-frame spike is treated differently from a sustained anomaly.
\end{itemize}

The four layers are combined into one fusion score by a fixed weighted sum, and the fused score is compared against four severity thresholds (LOW, MEDIUM, HIGH, CRITICAL). Only frames whose fused score crosses the configured alert threshold are turned into alerts, subject to a per-identifier debounce so that a single burst is not reported as dozens of near-duplicate alerts. The exact weight and threshold values are listed in Table~\ref{tab:app-fusion-config} (Appendix~\ref{app:config}).

\section{Hardware layer}

The hardware layer (\texttt{backend/hardware/can\_interface.py}) wraps the \texttt{python-can} library to support three transports: ELM327 over a serial/slcan connection, PEAK PCAN-USB, and native SocketCAN. It exposes connect, disconnect, and status endpoints, and feeds every received frame through the same \texttt{RealtimeEngine} instance used by replay, with the vehicle identifier selected in the Settings view. This design means the detection pipeline itself required no change to support live hardware once an adapter is available; only the transport layer differs, which is why the code-complete-but-hardware-blocked status described in Chapter~\ref{chap:requirements} was an acceptable risk to carry through the project.

\section{AI explanation engine}
\label{sec:design-ai}

Figure~\ref{fig:ai-workflow} shows the explanation workflow triggered when the analyst selects an alert.

\begin{figure}[H]
\centering
\begin{tikzpicture}[
    box/.style={draw, rounded corners, minimum width=3.4cm, minimum height=1cm, align=center, font=\small},
    >=Stealth
]
\node[box, fill=blue!8] (alert) {Selected alert\\(can\_id, score, reason)};
\node[box, fill=orange!8, right=1.2cm of alert] (ctx) {Context builder:\\current features +\\state stats + KB match +\\similar historical case};
\node[box, fill=green!8, below=1cm of ctx] (llm) {LLM provider chain:\\Gemini $\to$ Ollama};
\node[box, fill=yellow!12, below=1cm of alert] (rule) {Rule-based\\fallback template};
\node[box, fill=purple!10, right=1.2cm of llm] (out) {Summary, detail,\\recommendation,\\KB match label};

\draw[->] (alert) -- (ctx);
\draw[->] (ctx) -- (llm);
\draw[->] (llm) -- (out);
\draw[->] (llm.south) |- (rule.east) node[midway,below,font=\tiny]{if all providers fail};
\draw[->] (rule) -| (out.south);
\end{tikzpicture}
\caption{AI explanation workflow: context building, LLM provider chain, and rule-based fallback.}
\label{fig:ai-workflow}
\end{figure}

The context builder assembles four sources before calling the language model: the current alert's own features (identifier, frequency, entropy, fusion-layer scores), the expected statistics for the vehicle's current state (so the prompt can say, for example, that the observed frequency is far above the expected range), a keyword match against a small attack knowledge base, and the most similar historical anomaly found by cosine similarity against a persistent, growing store of previously explained anomalies. The assembled context is sent to a language-model provider chain — Gemini first, a local Ollama model second — and, if neither is reachable, a rule-based template still produces a complete, correctly parameterized explanation instead of failing the request, satisfying requirement NFR6.

\section{Class diagram (backend detection core)}

Figure~\ref{fig:class-diagram} shows the simplified class diagram of the backend's detection core, the part of the system most directly involved in the validation work of Chapter~\ref{chap:validation}.

\begin{figure}[H]
\centering
\begin{tikzpicture}[
    classbox/.style={draw, rectangle split, rectangle split parts=2, align=left, font=\scriptsize, text width=4.0cm},
    >=Stealth
]
\node[classbox] (engine) at (0,0) {
    \textbf{RealtimeEngine}
    \nodepart{two} + vehicle\_id: str\\ + process\_frame(frame)
};
\node[classbox] (scoring) at (5,0) {
    \textbf{ScoringEngine}
    \nodepart{two} + score(timing, payload, ml)\\ - weights, thresholds
};
\node[classbox] (store) at (0,-3.2) {
    \textbf{RuntimeModelStore}
    \nodepart{two} + classify\_state(features)\\ + score(state, features, vehicle)
};
\node[classbox] (bundle) at (5,-3.2) {
    \textbf{StateModelBundle}
    \nodepart{two} - models: dict[state, model]\\ + score(state, features)
};
\node[classbox] (stats) at (10,0) {
    \textbf{IDSStatistics}
    \nodepart{two} + on\_frame() / on\_alert()\\ + save\_reports()
};

\draw[->] (engine) -- (scoring);
\draw[->] (engine) -- (store);
\draw[->] (store) -- (bundle);
\draw[->] (engine) -- (stats);
\end{tikzpicture}
\caption{Simplified class diagram of the backend detection core.}
\label{fig:class-diagram}
\end{figure}

\texttt{RealtimeEngine} is the single entry point used by replay, live simulation, and hardware capture. It owns the current \texttt{vehicle\_id}, delegates timing/payload scoring and fusion to \texttt{ScoringEngine}, and delegates ML scoring to \texttt{RuntimeModelStore}, which holds one \texttt{StateModelBundle} per vehicle profile (plus a generic fallback bundle) and is responsible for selecting the correct per-state, per-vehicle model. \texttt{IDSStatistics} accumulates per-run counters and severity/alert-reason distributions and writes the final summary report consumed by both the desktop client and the \gls{ai} explanation engine.

\section{Deployment diagram}

Figure~\ref{fig:deployment} shows how the system is deployed on a single workstation, with the two external systems it may reach over the network.

\begin{figure}[H]
\centering
\begin{tikzpicture}[
    node/.style={draw, rounded corners, minimum width=5.6cm, minimum height=1cm, align=center, font=\small},
    ext/.style={draw, rounded corners, dashed, minimum width=3.6cm, minimum height=1cm, align=center, font=\small},
    >=Stealth
]
\node[draw, minimum width=8.4cm, minimum height=5.4cm] (host) at (0,0) {};
\node[font=\small\bfseries] at (-2.6,2.4) {Windows workstation};
\node[node, fill=blue!8] (wpf) at (0,1.4) {CANvisionNative.exe\\(WPF process)};
\node[node, fill=orange!8] (py) at (0,-0.6) {uvicorn / FastAPI\\(Python process, port 8765)};
\node[node, fill=gray!10, minimum height=0.8cm] (files) at (0,-2.2) {SQLite DB, model store,\\DBC files, log files};
\node[ext, fill=green!8] (llm) at (5.8,1.4) {LLM API\\(Gemini)};
\node[ext, fill=green!8] (ollama) at (5.8,-0.4) {Ollama\\(localhost, optional)};
\node[ext, fill=yellow!12] (can) at (-5.6,-0.6) {ELM327 / SocketCAN\\adapter + vehicle};

\draw[<->] (wpf) -- node[right,font=\scriptsize]{HTTP :8765} (py);
\draw[<->] (py) -- (files);
\draw[<->] (py) -- node[above,font=\scriptsize]{HTTPS} (llm);
\draw[<->] (py) -- node[above,font=\scriptsize]{localhost} (ollama);
\draw[<->] (py) -- node[above,font=\scriptsize]{serial / SocketCAN} (can);
\end{tikzpicture}
\caption{Deployment diagram: single-workstation deployment with external LLM and hardware links.}
\label{fig:deployment}
\end{figure}

\section{Conclusion}

This chapter presented the layered architecture of CANvisionNative, its desktop and backend components, its database schema, its \gls{rest} \gls{api}, the replay engine, the four-layer detection pipeline, the hardware transport layer, and the \gls{ai} explanation workflow, together with the class, entity-relationship, and deployment diagrams that summarize the design. The next chapter describes how this design was implemented.
